\documentclass[10pt,a4paper]{extarticle}
\pagenumbering{gobble}
\usepackage[utf8]{inputenc}
\usepackage[T1]{fontenc}
\usepackage[turkish]{babel}
\usepackage[left=0.6in,top=0.6in,right=0.6in,bottom=0.6in]{geometry}
\usepackage{hyperref}
\usepackage{titlesec}
\usepackage{enumitem}
\usepackage{datetime}
\usepackage{graphicx}
\titleformat{\section}{\large\bfseries\scshape\raggedright}{}{0em}{}[\titlerule]
\titlespacing*{\section}{0pt}{12pt}{8pt}
\begin{document}
\begin{center}
    \begin{minipage}{\textwidth}
        \centering
        {\LARGE\textbf{Yiğit Efe Avcı}} \hspace{2pt} {\LARGE{İstatistik ve Veri Bilimi}}\\[10pt]
        \href{mailto:yigit433.devs@gmail.com}{yigit433.devs@gmail.com} \textbullet\
        \href{https://github.com/yigit433}{github.com/yigit433} \textbullet\
        \href{https://linkedin.com/in/yigitefeavci}{linkedin.com/in/yigitefeavci} \textbullet\
        \href{https://yigit433.vercel.app}{yigit433.vercel.app}
    \end{minipage}
\end{center}
\section{Eğitim}
\textbf{İstanbul Ticaret Üniversitesi} \hfill İstanbul, Türkiye\\
\textit{İstatistik Lisans} \hfill 2023 -- 2027\\
\section{Profil}
Veri bilimi, istatistiksel modelleme ve anket araştırmalarına güçlü ilgi duyan istatistik lisans öğrencisiyim. Python, R ve SPSS başta olmak üzere modern veri kütüphaneleri (Pandas, Polars, Scikit-learn) ile uygulamalı deneyime sahibim. Yurt çapında saha araştırması deneyimimi; veri toplama, otomasyon ve zenginleştirme araçları geliştirmek için kullandığım yazılım mühendisliği becerileriyle birleştiriyorum.
\section{Deneyim}
\textbf{Konda Araştırma ve Danışmanlık} \hfill Türkiye\\
\textit{Saha Araştırmacısı (Anketör)} \hfill Ay Yıl -- Ay Yıl
\begin{itemize}[leftmargin=*,noitemsep,topsep=0pt]
    \item Yurt çapında yürütülen anketler için saha çalışmaları gerçekleştirdim; yüz yüze görüşmeler ve veri toplama süreçlerini kapsıyordu.
    \item Metodolojik standartlara uyarak ve katılımcılarla doğrudan etkileşim kurarak araştırma projelerine doğru ve güvenilir veri girdisi sağladım.
\end{itemize}
\textbf{WiFiCeps} \hfill Türkiye\\
\textit{Backend Developer (Proje Bazlı)} \hfill Ocak 2023 -- Şubat 2023
\begin{itemize}[leftmargin=*,noitemsep,topsep=0pt]
    \item Python tabanlı bir otomasyon çerçevesi ve FFmpeg destekli bir video işleme algoritması inşa ettim.
    \item Google Drive'a yüklenen videoları izleyen, güncellemeleri tespit edip platformla otomatik olarak senkronize eden bir sistem tasarladım.
    \item Otomatik işleme süreci ile video boyut ve kalitesini optimize ederek dosyaları yaklaşık \textasciitilde50MB seviyesine indirdim.
\end{itemize}
\textbf{\href{https://codeshare.me}{Code Share}} \hfill Londra, Birleşik Krallık\\
\textit{Yönetici Ortak (Serbest Meslek)} \hfill Ağustos 2020 -- Mayıs 2025\\
\textit{Yazılım Geliştirici (Freelance, Uzaktan)} \hfill Haziran 2019 -- Haziran 2023\\
\section{Projeler}
\textbf{\href{https://pickin.co}{Pickin} — Gerçek zamanlı çekiliş platformu} \hfill Next.js, TypeScript, PostgreSQL\\
\begin{itemize}[leftmargin=*,noitemsep,topsep=0pt]
    \item Redis Pub/Sub ile beslenen Server-Sent Events üzerinden anlık katılımcı ve durum güncellemeleri.
    \item NATS JetStream üzerinden zamanlanmış iş işleme ve OAuth tabanlı katılım şartları.
\end{itemize}
\textbf{\href{https://deckplane.com}{Deckplane} — Dağıtık Docker yönetim platformu} \hfill Bun, Hono, TypeScript, PostgreSQL\\
\begin{itemize}[leftmargin=*,noitemsep,topsep=0pt]
    \item Control-plane / host-agent mimarisi; ajanlar merkezi API ile WebSocket ve REST üzerinden haberleşir.
    \item GitHub repolarından dağıtım otomasyonu, RBAC, denetim kayıtları ve otomatik OpenAPI 3.0 dokümantasyonu.
\end{itemize}
\section{Yetenekler}
\textbf{Veri \& İstatistik:} Python, R, SPSS, Pandas, Polars, NumPy, Matplotlib, Seaborn, Scikit-learn, TensorFlow, PyTorch \\
\textbf{Diller:} Python, R, TypeScript, JavaScript, GoLang, Rust, Dart \\
\textbf{Mühendislik:} Node.js, Bun, Next.js, React, PostgreSQL, MongoDB, Docker \\
\textbf{Araçlar:} Jupyter, Git, Drizzle ORM, Postman
\end{document}
